\documentclass[10pt,a4paper]{amsart}
\usepackage{amsmath,amssymb,amsthm,graphicx}
% Retained preamble for the Stacks Project chapter snapshot in this directory.
% The chapter's own preamble.tex is a copy of the shared Stacks Project preamble; the block
% below is sliced MECHANICALLY, line for line, from it. Nothing was added, removed or
% redefined; the wrapper class, the xr-hyper/hyperref wiring and the externaldocument list
% are omitted because the assembler supplies the document class.
%
% lines taken from preamble.tex: 3, 5, 9, 11, 12, 156, 157, 158, ... (51 lines total)
%
% package coverage restored from the chapter's own preamble (2026-10-13): xy, xy-2cell

\usepackage[all]{xy}
\xyoption{2cell}
\UseAllTwocells

\newenvironment{reference}{\comment}{\endcomment}
\newenvironment{slogan}{\comment}{\endcomment}
\newenvironment{history}{\comment}{\endcomment}
\theoremstyle{plain}
\newtheorem{theorem}[subsection]{Theorem}
\newtheorem{proposition}[subsection]{Proposition}
\newtheorem{lemma}[subsection]{Lemma}
\theoremstyle{definition}
\newtheorem{definition}[subsection]{Definition}
\newtheorem{example}[subsection]{Example}
\newtheorem{exercise}[subsection]{Exercise}
\newtheorem{situation}[subsection]{Situation}
\theoremstyle{remark}
\newtheorem{remark}[subsection]{Remark}
\newtheorem{remarks}[subsection]{Remarks}
\numberwithin{equation}{subsection}
\def\lim{\mathop{\mathrm{lim}}\nolimits}
\def\colim{\mathop{\mathrm{colim}}\nolimits}
\def\Spec{\mathop{\mathrm{Spec}}}
\def\Hom{\mathop{\mathrm{Hom}}\nolimits}
\def\Ext{\mathop{\mathrm{Ext}}\nolimits}
\def\SheafHom{\mathop{\mathcal{H}\!\mathit{om}}\nolimits}
\def\SheafExt{\mathop{\mathcal{E}\!\mathit{xt}}\nolimits}
\def\Sch{\mathit{Sch}}
\def\Mor{\mathop{\mathrm{Mor}}\nolimits}
\def\Ob{\mathop{\mathrm{Ob}}\nolimits}
\def\Sh{\mathop{\mathit{Sh}}\nolimits}
\def\NL{\mathop{N\!L}\nolimits}
\def\CH{\mathop{\mathrm{CH}}\nolimits}
\def\proetale{{pro\text{-}\acute{e}tale}}
\def\etale{{\acute{e}tale}}
\def\QCoh{\mathit{QCoh}}
\def\Ker{\mathop{\mathrm{Ker}}}
\def\Im{\mathop{\mathrm{Im}}}
\def\Coker{\mathop{\mathrm{Coker}}}
\def\Coim{\mathop{\mathrm{Coim}}}
\DeclareMathSymbol{\boxtimes}{\mathbin}{AMSa}{"02}
\def\QCohstack{\mathcal{QC}\!\mathit{oh}}
\def\Cohstack{\mathcal{C}\!\mathit{oh}}
\def\Spacesstack{\mathcal{S}\!\mathit{paces}}
\def\Quotfunctor{\mathrm{Quot}}
\def\Hilbfunctor{\mathrm{Hilb}}
\def\Curvesstack{\mathcal{C}\!\mathit{urves}}
\def\Polarizedstack{\mathcal{P}\!\mathit{olarized}}
\def\Complexesstack{\mathcal{C}\!\mathit{omplexes}}
\def\Pic{\mathop{\mathrm{Pic}}\nolimits}
\def\Picardstack{\mathcal{P}\!\mathit{ic}}
\def\Picardfunctor{\mathrm{Pic}}
\def\Deformationcategory{\mathcal{D}\!\mathit{ef}}
\usepackage{cleveref}
\crefname{originallemma}{original lemma}{original lemmas}
\Crefname{originallemma}{Original lemma}{Original lemmas}
\makeatletter
\newlabel{example-universal-quasi-finite-syntomic}{{0FK8}{source}{}{sourcepointer-1}{}}
\newlabel{example-universal-quasi-finite-syntomic@cref}{{[originallemma][0][]0FK8}{[1][0][]source}{}{}{}}
\makeatother
\crefname{originallemma}{original lemma}{original lemmas}
\Crefname{originallemma}{Original lemma}{Original lemmas}
\makeatletter
\newlabel{lemma-syntomic-quasi-finite}{{0BWE}{source}{}{sourcepointer-2}{}}
\newlabel{lemma-syntomic-quasi-finite@cref}{{[originallemma][0][]0BWE}{[1][0][]source}{}{}{}}
\makeatother
\crefname{originallemma}{original lemma}{original lemmas}
\Crefname{originallemma}{Original lemma}{Original lemmas}
\makeatletter
\newlabel{morphisms-lemma-etale-at-point}{{02GU}{source}{}{sourcepointer-3}{}}
\newlabel{morphisms-lemma-etale-at-point@cref}{{[originallemma][0][]02GU}{[1][0][]source}{}{}{}}
\makeatother
\title{Stacks, Discriminants: lemma-universal-quasi-finite-syntomic-etale}
\author{The Stacks Project Authors}
\date{stacks-project:discriminant.tex:local-snapshot-20261012; GNU FDL 1.2 or later}
\begin{document}
\maketitle
\noindent\textit{Editorial note.} Statement span is lines 2187-2195 of sources/stacks-discriminant/pinned-original-source.tex; proof span is lines 2197-2222, the complete author proof (24 body lines).\par\medskip
\noindent\textit{Editorial note.} Navigation stubs: 3 cited-but-not-included labels are supplied as non-load-bearing pointers with their pinned tags where the snapshot tag index has them.\par\medskip
\noindent\textit{Editorial note.} extra\_packages is empty on purpose: the retained preamble already loads xy with options and a second declaration clashes in the assembler wrapper.\par\medskip
\noindent\textit{Editorial note.} Text is retained verbatim, including the author's own typos and placeholders.\par\medskip
\noindent\textit{Editorial note.} Queue index 23 of queue\_discriminant.json (pb=1174, stmt 2187-2195, proof 2197-2222).\par\medskip
\noindent\textit{Allowed standard prerequisites (editorial).}\begin{itemize}
\item Explicit prerequisite boundary. The cited inputs are prerequisites; the argument in the retained span is not replaced by any cited result.
\end{itemize}\medskip
\ifdefined\hypertarget\hypertarget{sourcepointer-1}{}\fi
\noindent\textit{Source reference.} Original lemma 0FK8: sources/stacks-discriminant/pinned-original-source.tex, source lines 2168--2185. Cited by the retained statement or proof but not part of the collected spans; supplied only as a navigation pointer. The argument does not depend on its content.\par\medskip
\ifdefined\hypertarget\hypertarget{sourcepointer-2}{}\fi
\noindent\textit{Source reference.} Original lemma 0BWE: sources/stacks-discriminant/pinned-original-source.tex, source lines 2005--2026. Cited by the retained statement or proof but not part of the collected spans; supplied only as a navigation pointer. The argument does not depend on its content.\par\medskip
\ifdefined\hypertarget\hypertarget{sourcepointer-3}{}\fi
\noindent\textit{Source reference.} Original lemma 02GU: sources/stacks-discriminant/pinned-original-source.tex, source lines 2222--2222. Cited by the retained statement or proof but not part of the collected spans; supplied only as a navigation pointer. The argument does not depend on its content.\par\medskip
\begin{lemma}
\label{lemma-universal-quasi-finite-syntomic-etale}
With notation as in Example \ref{example-universal-quasi-finite-syntomic}
the schemes $X_{n, d}$ and $Y_{n, d}$ are regular and irreducible,
the morphism $Y_{n, d} \to X_{n, d}$ is locally quasi-finite and
syntomic, and there is a dense open subscheme $V \subset Y_{n, d}$
such that $Y_{n, d} \to X_{n, d}$ restricts to an \'etale morphism
$V \to X_{n, d}$.
\end{lemma}

\begin{proof}
The scheme $X_{n, d}$ is the spectrum of the polynomial ring $A$.
Hence $X_{n, d}$ is regular and irreducible. Since we can write
$$
f_i = a_{i, (0, \ldots, 0)} +
\sum\nolimits_{E \in T, E \not = (0, \ldots, 0)} a_{i, E} x^E
$$
we see that the ring $B$ is isomorphic to the polynomial ring
on $x_1, \ldots, x_n$ and the elements $a_{i, E}$ with
$E \not = (0, \ldots, 0)$. Hence $\Spec(B)$ is an irreducible and
regular scheme and so is the open $Y_{n, d}$. The morphism
$Y_{n, d} \to X_{n, d}$ is locally quasi-finite and syntomic by
Lemma \ref{lemma-syntomic-quasi-finite}. To find $V$ it suffices
to find a single point where $Y_{n, d} \to X_{n, d}$ is \'etale
(the locus of points where a morphism is \'etale is open by
definition). Thus it suffices to find a point of $X_{n, d}$
where the fibre of $Y_{n, d} \to X_{n, d}$ is nonempty and \'etale, see
Morphisms, Lemma \ref{morphisms-lemma-etale-at-point}. We choose
the point corresponding to the ring map $\chi : A \to \mathbf{Q}$
sending $f_i$ to $x_i^d - 1$. Then
$$
B \otimes_{A, \chi} \mathbf{Q} =
\mathbf{Q}[x_1, \ldots, x_n]/(x_1^d - 1, \ldots, x_n^d - 1)
$$
which is a nonzero \'etale algebra over $\mathbf{Q}$.
\end{proof}

\end{document}
